\chapter{Introduction}\label{ch:intro}

\begin{chapterabstract}
This chapter introduces the $k$-dispersion problem, explains why points in
convex position are an interesting special case, states the problem
precisely, lists the objectives of the project and summarises its
contributions.
\end{chapterabstract}

\section{Background}

Facility location asks where to place a number of facilities so that some
objective is optimised. In the classical versions, such as $k$-center or
$k$-median, facilities should be \emph{close} to the clients they serve. In
\emph{obnoxious} facility location the opposite is true: facilities should be
far from one another, because they interfere with or endanger each other.
Typical examples are radio transmitters that must not interfere, franchise
outlets that should not compete for the same customers, and storage plants for
hazardous material that should be placed so that an accident at one plant does
not reach another~\citep{tamir,ravi}.

The most basic problem of this kind is \emph{max--min dispersion}, also called
$k$-dispersion or the maximally separated set problem~\citep{agarwal}: given a
set $P$ of $n$ candidate locations and an integer $k$, choose $k$ of them so
that the smallest distance between two chosen locations is as large as
possible. For points in the plane with Euclidean distances the problem is
NP-hard~\citep{wangkuo}; the best exact algorithms take $n^{O(\sqrt k)}$
time~\citep{akagi}, and a ratio of $1/2$ is the best known approximation
guarantee in polynomial time~\citep{ravi,tamir}.

A natural way to obtain efficient algorithms is to restrict the input. When the
candidate locations lie on the boundary of a convex region---for example, on
the coast of a convex island, the setting used to motivate the problem
in~\citet{sbm}---the points are in \emph{convex position}, and the problem
becomes solvable in polynomial time. This project studies $k$-dispersion for
points in convex position. Our starting point is the recent paper of
Singireddy, Basappa and Mitchell in \emph{Discrete Applied
Mathematics}~\citep{sbm}, which gives a fixed-parameter algorithm based on a
bounded search tree, a polynomial dynamic program and an approximation for
$k=3$.

\section{Problem statement}\label{sec:problem}

Let $P$ be a set of $n$ points in the plane in \emph{convex position}, that is,
every point of $P$ is a vertex of the convex hull $\conv(P)$. For an integer
$2\le k\le n$, the \emph{$k$-dispersion} problem asks for a subset
$S\subseteq P$ with $|S|=k$ that maximises
\[
  \delta(S)=\min_{p\ne q\in S}|pq| .
\]
We write $\rstar=\rstar_k(P)$ for the optimum. The \emph{decision version}
asks, for a given radius $r$, whether $P$ contains $k$ points whose pairwise
distances are all at least $r$ (an \emph{$r$-far $k$-set}). Since $\rstar$ is
one of the $\binom n2$ pairwise distances and the decision is monotone in $r$,
an algorithm for the decision version yields one for the optimisation version
by binary search.

\section{Motivation}

Two questions motivated the project.

\begin{enumerate}[leftmargin=*,itemsep=2pt]
\item \emph{Is the search tree of~\citet{sbm} correct?} Its correctness rests on
a claim (Claim~2 of~\citet{sbm}) that, if $k$ well-separated points exist, a
solution can be grown from some starting point by repeatedly adding the first
admissible point at one of the two ends of the chosen arc. The published proof
of this claim is a picture argument. A rigorous proof would put the
algorithm on solid ground; a counterexample would show that it is incorrect.
\item \emph{When is the problem easy?} For $k=3$ much faster algorithms are
known than for general $k$~\citep{kobayashi,tw25}. We wanted to understand
which structural properties make $k=3$ easy, whether they extend to larger
$k$, and for which inputs a simple greedy rule is exact.
\end{enumerate}

\section{Objectives of the project}

\begin{itemize}[leftmargin=*,itemsep=2pt]
\item Study the algorithms of~\citet{sbm} and the related work, and re-implement
them in exact integer arithmetic.
\item Prove or refute the claim on which the bounded search tree of~\citet{sbm}
relies, at least for small $k$.
\item Determine which greedy properties that hold for $k=3$ extend to larger
$k$, and identify a natural class of inputs on which greedy is exact.
\item Obtain a faster algorithm for such a special class, and a simple exact
algorithm for the general problem.
\item Implement all algorithms, validate them against exhaustive search, and
evaluate them experimentally.
\end{itemize}

\section{Contributions}\label{sec:contributions}

\begin{enumerate}[leftmargin=*,itemsep=2pt]
\item \textbf{The two-sided search tree} (\cref{sec:sbm}). We prove the claim
of~\citet{sbm} for $k\le4$ (\cref{thm:sbm4}), which makes their algorithm
correct for these $k$. For every $k$ we prove an attachment lemma that
certifies single steps of the search and show, using Delaunay ears, that the
first step can always be made safely (\cref{thm:firststep}). We also give point
sets showing that the starting point cannot be chosen by distances or angles
alone (\cref{prop:rootrules}).

\item \textbf{Greedy structure} (\cref{sec:greedy}). For $k=3$ every $r$-far
triple can be turned into one of a few canonical forms by a single
replacement, and an arc characterisation of $r$-far triples holds. For
$k\ge4$ these properties fail; in particular there are point sets in which no
solution contains a point together with the first point at distance at least
$r$ from it (\cref{thm:tightpair}). The $8$-point example for $k=4$ can be
checked by hand.

\item \textbf{The arc property} (\cref{sec:arc}). If, for every point $q$, the
points at distance at least $r$ from $q$ are contiguous along the hull, then a
single greedy chain decides the problem (\cref{thm:arc}). For unimodal point
sets, which include perturbed points on a circle, this gives an exact
algorithm with expected running time $O(n\log^2n)$ (\cref{thm:arcopt}).

\item \textbf{A simple exact algorithm} (\cref{ch:exact}). A set in convex
position is $r$-far if and only if it has a triangulation with long edges and
locally Delaunay diagonals (\cref{thm:cert-delaunay}). This certificate gives a
short correctness proof for a dynamic program over chords; capping the
per-chord information at $k-2$ thresholds decides the problem in
$O(n^3\log k)$ time and solves it in $O(n^3\log k\log n)$ time.

\item \textbf{Implementation and evaluation} (\cref{ch:results}). All
algorithms are implemented in C++ with exact integer predicates, cross-checked
against exhaustive search and independent code, and evaluated on instances
with up to $10^6$ points.
\end{enumerate}

\section{Organisation of the report}

\textbf{Chapter~\ref{ch:lit}} reviews the literature.
\textbf{Chapter~\ref{ch:prelim}} fixes notation and proves the Delaunay
certificate. \textbf{Chapter~\ref{sec:greedy}} studies greedy structure for
$k=3$ and $k\ge4$. \textbf{Chapter~\ref{sec:sbm}} analyses the two-sided search
tree of~\citet{sbm} and proves its correctness for $k\le4$.
\textbf{Chapter~\ref{sec:arc}} treats point sets with the arc property.
\textbf{Chapter~\ref{ch:exact}} presents the exact algorithm.
\textbf{Chapter~\ref{ch:results}} describes the implementation and the
experimental results, and \textbf{Chapter~\ref{ch:conclusion}} concludes with
limitations and future work.
